% Volume IV: Optimization and Computation
% Continuity note for Chapter 23.
% Previous dependencies: Chapters 4-7 provide matrices as typed linear maps, tensor shapes, Frobenius pairings, trace tricks, Jacobians, Hessians, and chain rules; Chapters 8-10 provide derivatives as local linear maps and Taylor approximation; Chapters 20-22 provide losses, gradients, Newton ideas, constraints, and variational objectives.
% New durable concepts: computational graphs; forward evaluation; topological order; local derivatives; Jacobian-vector products; vector-Jacobian products; forward-mode automatic differentiation; reverse-mode automatic differentiation; adjoints; shape discipline for matrix and tensor gradients; floating-point error; overflow, underflow, cancellation; log-sum-exp; stable softmax and cross-entropy; vectorization, broadcasting, batching, BLAS/GPU kernels, memory movement, and efficient scientific computation.
% Forward hooks: Chapter 24 studies stable numerical linear algebra; Chapter 38 develops backpropagation and differentiable programming; Chapters 39-42 use autodiff in deep optimization and generative modeling; Chapters 44-51 use differentiable simulations, neural data likelihoods, and model fitting.
% Notation commitments: Computational graph nodes are v_i; local derivatives are partial v_j/partial v_i along graph edges; tangent variables use dot notation \dot v_i; reverse adjoints use \bar v_i=\partial L/\partial v_i for scalar losses; Jacobians use numerator layout J_{ij}=\partial y_i/\partial x_j; scalar-loss matrix gradients satisfy dL=\langle \nabla_W L,dW\rangle_F=\operatorname{tr}((\nabla_W L)^\top dW).
\section{Chapter 23: Automatic Differentiation and Numerical Computation}
\label{ch:automatic-differentiation-numerical-computation}

The optimization chapters taught us to write losses and gradients. Chapter 7 taught us that gradients, Jacobians, and Hessians have shapes, not just formulas. Modern scientific computing asks how those derivatives are actually computed inside a program that may contain thousands of array operations, branches, matrix multiplications, nonlinearities, and reductions.

The central idea is simple but powerful. A program is decomposed into elementary operations. Each operation has a local derivative. Automatic differentiation applies the chain rule through the program exactly up to floating-point arithmetic. This is different from symbolic differentiation, which manipulates expressions, and different from finite differences, which approximate derivatives by rerunning the program at perturbed inputs.

\paragraph{Chapter problem.}
How can a numerical program compute correct, shape-consistent derivatives while also respecting finite precision, stability, batching, and efficient hardware execution?

\paragraph{Section map.}
Section 23.1 defines computational graphs and forward evaluation. Section 23.2 derives forward-mode and reverse-mode automatic differentiation through Jacobian-vector products and vector-Jacobian products. Section 23.3 explains floating-point precision, numerical stability, stable softmax, and log-sum-exp. Section 23.4 develops vectorization, batching, broadcasting, and efficient array computation.

\subsection{23.1 Computational graphs}

\paragraph{Guiding question.}
How can a complicated formula or program be represented as elementary operations so that the chain rule can be applied locally?

\paragraph{Beginner intuition.}
A computational graph is a bookkeeping device for a calculation. Each node stores an intermediate value, and each directed edge says that one value is used to compute another. The graph matters because derivatives also move along these edges. If a node affects the loss through several downstream paths, its total derivative contribution is the sum of the contributions from all paths.

For example, define
\[
z=(xy+\sin x)^2.
\]
Instead of differentiating this expression in one step, introduce intermediate variables
\[
v_1=x,\qquad v_2=y,\qquad v_3=v_1v_2,\qquad
v_4=\sin v_1,\qquad v_5=v_3+v_4,\qquad v_6=v_5^2.
\]
The output is \(z=v_6\). Each line is elementary: copy, multiply, sine, add, square. This makes the chain rule mechanical.

\paragraph{Formal setup.}
A finite directed acyclic computational graph consists of nodes \(v_1,\ldots,v_m\) arranged in a topological order: whenever \(v_j\) depends on \(v_i\), we have \(i<j\). Input nodes are assigned external values. Each non-input node is computed by a local function
\[
v_j=\phi_j(v_{p_1},\ldots,v_{p_r}),
\]
where \(p_1,\ldots,p_r<j\) are parent nodes. A forward pass evaluates the nodes in topological order.

If the graph output is a scalar loss \(L=v_m\), then the derivative of \(L\) with respect to an input is obtained by summing over all directed paths from that input to the output. For one path
\[
v_{i_0}\to v_{i_1}\to\cdots\to v_{i_k}=L,
\]
the path contribution is
\[
\frac{\partial v_{i_1}}{\partial v_{i_0}}
\frac{\partial v_{i_2}}{\partial v_{i_1}}
\cdots
\frac{\partial L}{\partial v_{i_{k-1}}},
\]
with multiplication intended between the factors. The total derivative sums these products over all paths. Reverse-mode autodiff will compute this sum without explicitly listing the paths.

\paragraph{Worked mechanism.}
Use \(z=(xy+\sin x)^2\) at \(x=2\), \(y=3\). The forward values are
\[
v_1=2,\quad v_2=3,\quad v_3=6,\quad v_4=\sin 2,\quad
v_5=6+\sin 2,\quad v_6=(6+\sin 2)^2.
\]
The local derivatives are
\[
\frac{\partial v_3}{\partial v_1}=v_2,\qquad
\frac{\partial v_3}{\partial v_2}=v_1,\qquad
\frac{\partial v_4}{\partial v_1}=\cos v_1,
\]
\[
\frac{\partial v_5}{\partial v_3}=1,\qquad
\frac{\partial v_5}{\partial v_4}=1,\qquad
\frac{\partial v_6}{\partial v_5}=2v_5.
\]
There are two paths from \(x=v_1\) to \(z=v_6\): through the product node \(v_3\) and through the sine node \(v_4\). Therefore
\[
\frac{\partial z}{\partial x}
=2v_5\left(v_2+\cos v_1\right)
=2(6+\sin 2)(3+\cos 2).
\]
There is one path from \(y=v_2\) to \(z\):
\[
\frac{\partial z}{\partial y}
=2v_5v_1
=4(6+\sin 2).
\]
This is ordinary calculus, but the graph forces us to account for every route by which a variable affects the output.

\paragraph{Numerical, neuroscience, and AI examples.}
A numerical example is the loss
\[
L(a,b)=\frac12(a b-1)^2.
\]
With \(a=2\), \(b=0.4\), the intermediate prediction is \(ab=0.8\), the residual is \(-0.2\), and \(L=0.02\). The graph shows that \(\partial L/\partial a=(ab-1)b=-0.08\) and \(\partial L/\partial b=(ab-1)a=-0.4\).

In computational neuroscience, a differentiable firing-rate model might compute
\[
r=\sigma(wx+b),\qquad L=\frac12(r-r_{\mathrm{obs}})^2.
\]
The graph nodes are stimulus \(x\), synaptic weight \(w\), drive \(u=wx+b\), rate \(r=\sigma(u)\), residual, and loss. Local derivatives identify how a change in \(w\) changes the predicted rate and therefore the data-fitting loss.

In AI, a neural network is a large computational graph. A transformer block contains matrix multiplications, softmax operations, residual additions, and normalizations. The graph view says that training is not magic: each operation supplies local derivative rules, and the loss gradient is composed through the graph.

\paragraph{Common misconceptions and failure modes.}
A graph is not the same as a drawing of the final formula. It records intermediate values and dependency order. A variable reused in two places creates two downstream paths, and both contribute to the derivative. Another common failure is to forget shapes: if \(v_j\) is a vector, its local derivative with respect to a vector parent is a Jacobian, not a scalar. Finally, a computational graph for a program is not always acyclic if the program contains loops, but each finite unrolling of the loop is an acyclic graph.

\paragraph{Exercises.}
\begin{enumerate}[label=\textbf{23.1.\arabic*.}, leftmargin=*]
\item \textbf{Mechanical.} Build a computational graph for \(L=(x^2+y)^3\). List the nodes in topological order.
\item \textbf{Proof or derivation.} For \(z=(xy+\sin x)^2\), derive \(\partial z/\partial y\) by multiplying local derivatives along graph paths.
\item \textbf{Numerical/computational.} Evaluate the forward graph for \(L=\frac12(ab-1)^2\) at \(a=1.5\), \(b=0.8\), and compute the local derivative at each non-input node.
\item \textbf{Interpretation.} In the firing-rate graph \(r=\sigma(wx+b)\), \(L=\frac12(r-r_{\mathrm{obs}})^2\), explain what each node represents experimentally.
\item \textbf{Misconception/debugging.} A student differentiates \(z=(xy+\sin x)^2\) with respect to \(x\) but includes only the path through \(xy\). Identify the missing path and correct the derivative.
\end{enumerate}

\subsection{23.2 Forward-mode and reverse-mode autodiff}

\paragraph{Guiding question.}
Once a computational graph is built, should derivatives be propagated forward from inputs or backward from a scalar loss?

\paragraph{Beginner intuition.}
Forward mode asks: if the input is nudged in a particular direction, how do all later values change? Reverse mode asks: if the output loss changes, how sensitive is it to each earlier value? Both are exact applications of the chain rule on the same graph. They differ in what they compute efficiently.

For a function \(f:\mathbb{R}^n\to\mathbb{R}^m\), the Jacobian \(J_f(x)\in\mathbb{R}^{m\times n}\) may be too large to form. Forward mode naturally computes products \(J_f(x)u\), called Jacobian-vector products. Reverse mode naturally computes products \(w^\top J_f(x)\), called vector-Jacobian products. For a scalar loss \(L:\mathbb{R}^n\to\mathbb{R}\), reverse mode computes the full gradient \(\nabla_x L\) in one backward sweep, which is why backpropagation is efficient for neural networks with many parameters.

\paragraph{Formal setup.}
For each node \(v_j\), forward mode stores a primal value \(v_j\) and a tangent value \(\dot v_j\), where
\[
\dot v_j=\frac{d}{d\epsilon}v_j(x+\epsilon u)\bigg|_{\epsilon=0}
\]
for a chosen input direction \(u\). If
\[
v_j=\phi_j(v_{p_1},\ldots,v_{p_r}),
\]
then
\[
\dot v_j
=\sum_{\ell=1}^r
\frac{\partial \phi_j}{\partial v_{p_\ell}}\dot v_{p_\ell}.
\]
For vector-valued nodes, the local partial derivative is a Jacobian block and the formula is the same linear map applied to tangent vectors.

For a scalar loss \(L=v_m\), reverse mode stores adjoints
\[
\bar v_i=\nabla_{v_i} L,
\]
as column vectors with the same shape as \(v_i\). For scalar nodes this is the
ordinary partial derivative \(\partial L/\partial v_i\).
Initialize
\[
\bar v_m=1.
\]
Then traverse the graph backward. If \(v_i\) is a parent of \(v_j\), add the contribution
\[
\bar v_i \leftarrow \bar v_i+
\left(\frac{\partial v_j}{\partial v_i}\right)^\top\bar v_j.
\]
If \(v_i\in\mathbb{R}^{d_i}\) and \(v_j\in\mathbb{R}^{d_j}\), then
\(\partial v_j/\partial v_i\in\mathbb{R}^{d_j\times d_i}\), so the transpose
maps \(\bar v_j\in\mathbb{R}^{d_j}\) back to
\(\bar v_i\in\mathbb{R}^{d_i}\). The scalar formula
\(\bar v_i\leftarrow \bar v_i+\bar v_j(\partial v_j/\partial v_i)\) is the
special case \(d_i=d_j=1\). This is the column-adjoint version of a
vector-Jacobian product.

\paragraph{Worked mechanism.}
Return to
\[
v_3=v_1v_2,\quad v_4=\sin v_1,\quad v_5=v_3+v_4,\quad v_6=v_5^2.
\]
To compute \(\partial v_6/\partial x\) by forward mode, seed \(\dot v_1=1\), \(\dot v_2=0\). Then
\[
\dot v_3=v_2\dot v_1+v_1\dot v_2=v_2,\qquad
\dot v_4=\cos(v_1)\dot v_1=\cos(v_1),
\]
\[
\dot v_5=\dot v_3+\dot v_4=v_2+\cos(v_1),\qquad
\dot v_6=2v_5\dot v_5.
\]
Thus \(\dot v_6=\partial z/\partial x\).

Reverse mode initializes \(\bar v_6=1\). Then
\[
\bar v_5=2v_5,\qquad
\bar v_3=\bar v_5,\qquad
\bar v_4=\bar v_5,
\]
\[
\bar v_1=\bar v_3v_2+\bar v_4\cos(v_1),\qquad
\bar v_2=\bar v_3v_1.
\]
So
\[
\bar v_1=2v_5(v_2+\cos v_1),\qquad
\bar v_2=2v_5v_1.
\]
One reverse sweep gives both input derivatives.

\paragraph{Matrix and tensor gradients.}
Let
\[
y=Wx+b,\qquad
W\in\mathbb{R}^{m\times n},\quad x\in\mathbb{R}^n,\quad b,y\in\mathbb{R}^m,
\]
and let \(g=\nabla_y L\in\mathbb{R}^m\) be the incoming reverse-mode adjoint. The differential is
\[
dy=dW\,x+W\,dx+db.
\]
Because
\[
dL=g^\top dy,
\]
we get
\[
dL=g^\top dW\,x+g^\top W\,dx+g^\top db.
\]
Using the Frobenius pairing,
\[
g^\top dW\,x
=\operatorname{tr}(xg^\top dW)
=\operatorname{tr}((gx^\top)^\top dW).
\]
Therefore
\[
\boxed{\nabla_W L=gx^\top,\qquad \nabla_x L=W^\top g,\qquad \nabla_b L=g.}
\]
The shapes are
\[
\nabla_W L\in\mathbb{R}^{m\times n},\quad
\nabla_x L\in\mathbb{R}^n,\quad
\nabla_b L\in\mathbb{R}^m.
\]
This is the local matrix derivative used repeatedly in backpropagation.

For least squares,
\[
L(x)=\frac12\|Ax-b\|^2,
\]
the graph \(x\mapsto r=Ax-b\mapsto L=\frac12 r^\top r\) gives
\[
\nabla_x L=A^\top(Ax-b),\qquad
\nabla_x^2L=A^\top A.
\]
The Hessian is a shaped linear map from perturbations \(dx\in\mathbb{R}^n\) to gradient perturbations. For a quadratic form \(q(x)=x^\top Hx\) with \(H=H^\top\),
\[
dq=2(Hx)^\top dx,
\qquad
\nabla_x q=2Hx.
\]

\paragraph{Examples and connections.}
Numerically, if \(W=\begin{bmatrix}1&2\\3&4\end{bmatrix}\), \(x=(5,6)^\top\), and \(g=(0.1,-0.2)^\top\), then
\[
\nabla_W L=gx^\top=
\begin{bmatrix}
0.5&0.6\\
-1.0&-1.2
\end{bmatrix},
\qquad
\nabla_x L=W^\top g=(-0.5,-0.6)^\top.
\]
In neuroscience, \(W\) may be a linear decoder from neural population activity \(x\) to a latent behavioral variable \(y\). The formula \(W^\top g\) pulls a behavioral error sensitivity back to neural activity coordinates. In AI, the same calculation is the backward pass for a dense layer.

\paragraph{Common misconceptions and failure modes.}
Forward mode is not less correct than reverse mode; it is efficient for different shapes of problems. It is good when there are few inputs or when a directional derivative is needed. Reverse mode is good for scalar losses with many parameters. A second misconception is that reverse mode forms full Jacobians. Usually it does not; it composes vector-Jacobian products. A third failure is to ignore memory: reverse mode needs forward intermediate values, so large graphs may require checkpointing or recomputation.

\paragraph{Exercises.}
\begin{enumerate}[label=\textbf{23.2.\arabic*.}, leftmargin=*]
\item \textbf{Mechanical.} For \(f(x,y)=x^2y+\exp y\), define a computational graph and write the forward-mode tangent updates for the seed \((\dot x,\dot y)=(1,0)\).
\item \textbf{Proof or derivation.} Starting from \(dy=dW\,x+W\,dx+db\) and \(dL=g^\top dy\), derive \(\nabla_W L=gx^\top\) using the Frobenius pairing.
\item \textbf{Numerical/computational.} Let \(A=\begin{bmatrix}1&0\\2&1\\0&1\end{bmatrix}\), \(x=(1,2)^\top\), and \(b=(1,1,0)^\top\). Compute \(\nabla_x \frac12\|Ax-b\|^2\).
\item \textbf{Interpretation.} Explain why reverse mode is well matched to training a neural network with millions of parameters and one scalar minibatch loss.
\item \textbf{Misconception/debugging.} A student says, "Backpropagation stores the full Jacobian of every layer." Explain what is usually stored and what products are computed instead.
\end{enumerate}

\subsection{23.3 Numerical precision and stability}

\paragraph{Guiding question.}
Why can mathematically equivalent formulas behave differently on a computer?

\paragraph{Beginner intuition.}
Real numbers in mathematics have infinitely many digits. Floating-point numbers in a computer do not. They store a finite approximation with limited range and limited precision. As a result, arithmetic can round, overflow, underflow, and lose significant digits through cancellation. Numerical stability asks whether an algorithm controls these errors rather than amplifying them unnecessarily.

\paragraph{Formal setup.}
A simple floating-point model writes one elementary operation as
\[
\operatorname{fl}(a\circ b)=(a\circ b)(1+\delta),
\qquad |\delta|\leq \varepsilon_{\mathrm{mach}},
\]
where \(\circ\) is \(+,-,\times,\) or \(/\) and \(\varepsilon_{\mathrm{mach}}\) is machine epsilon for the chosen floating-point format. This model is an idealization, but it captures the main point: each operation is rounded.

A problem is \emph{ill-conditioned} if small input perturbations can cause large output perturbations. An algorithm is \emph{unstable} if it introduces large error for a problem that could be solved more accurately. Conditioning belongs to the mathematical problem; stability belongs to the computational method.

Three common failure modes are:
\[
\text{overflow: a magnitude exceeds representable range,}
\]
\[
\text{underflow: a positive magnitude rounds to zero or subnormal scale,}
\]
\[
\text{cancellation: subtracting nearly equal numbers loses leading digits.}
\]

\paragraph{Worked mechanism: log-sum-exp and stable softmax.}
The softmax of scores \(s\in\mathbb{R}^K\) is
\[
p_i=\frac{\exp(s_i)}{\sum_{j=1}^K\exp(s_j)}.
\]
If some \(s_i\) is large, \(\exp(s_i)\) may overflow. The mathematically equivalent stable form subtracts
\[
m=\max_j s_j.
\]
Then
\[
p_i=\frac{\exp(s_i-m)}{\sum_{j=1}^K\exp(s_j-m)}.
\]
All exponents \(s_i-m\leq 0\), so the largest exponential is \(1\).

The same idea stabilizes
\[
\log\sum_{j=1}^K\exp(s_j).
\]
Because
\[
\sum_j\exp(s_j)=\exp(m)\sum_j\exp(s_j-m),
\]
we get
\[
\boxed{\log\sum_j\exp(s_j)
=m+\log\sum_j\exp(s_j-m).}
\]
This identity is exact in real arithmetic and much safer in floating-point arithmetic.

For one-hot target \(y\in\mathbb{R}^K\) and softmax probabilities \(p\), the cross-entropy loss is
\[
L(s,y)=-\sum_{i=1}^K y_i\log p_i.
\]
Using \(p_i=\exp(s_i)/\sum_j\exp(s_j)\),
\[
L(s,y)=\log\sum_j\exp(s_j)-\sum_i y_i s_i.
\]
Differentiating gives
\[
\boxed{\nabla_s L=p-y.}
\]
This derivative is simple, but the probability vector \(p\) should be computed with the stable softmax.

\paragraph{Examples and connections.}
Numerically, the naive expression \(\log(\exp(1000)+\exp(999))\) overflows in many formats. The stable version uses \(m=1000\):
\[
1000+\log(1+\exp(-1)).
\]
This is finite and accurate.

Cancellation appears in
\[
\sqrt{x+1}-\sqrt{x}
\]
for large \(x\). Multiplying by the conjugate gives the equivalent expression
\[
\frac{1}{\sqrt{x+1}+\sqrt{x}},
\]
which avoids subtracting nearly equal large numbers.

In computational neuroscience, likelihoods for long spike trains can be products of many small probabilities. Direct products underflow, so models use log-likelihoods and sums. In AI, stable softmax and cross-entropy are essential for classification, attention, and language modeling. Numerical mistakes here can produce NaNs, zero gradients, or misleading confidence.

\paragraph{Common misconceptions and failure modes.}
Mathematical equality does not imply equal numerical behavior. A second misconception is that using more precision fixes every issue; it often helps but does not remove bad conditioning or unstable algorithms. A third failure is to compute probabilities first and then take logs when the log-domain calculation is available. Finally, finite-difference gradient checks can fail if the step size is too large, causing truncation error, or too small, causing roundoff error.

\paragraph{Exercises.}
\begin{enumerate}[label=\textbf{23.3.\arabic*.}, leftmargin=*]
\item \textbf{Mechanical.} Rewrite \(\log(\exp(4)+\exp(7)+\exp(6))\) using the log-sum-exp stabilization with \(m=\max_i s_i\).
\item \textbf{Proof or derivation.} Derive \(\nabla_s L=p-y\) for softmax probabilities and one-hot cross-entropy.
\item \textbf{Numerical/computational.} Compute the stable softmax for scores \((1000,999,998)\) up to expressions involving \(\exp(-1)\) and \(\exp(-2)\). Why is the naive computation dangerous?
\item \textbf{Interpretation.} In a Poisson spike-train likelihood with many small factors, why is it safer to maximize a log-likelihood than a raw likelihood product?
\item \textbf{Misconception/debugging.} A training run returns NaN after a softmax layer with very large logits. Identify the likely numerical failure and state the stable replacement.
\end{enumerate}

\subsection{23.4 Vectorization and efficient computation}

\paragraph{Guiding question.}
How do we express computations so that their mathematical shape, memory layout, and hardware execution are all efficient?

\paragraph{Beginner intuition.}
Vectorization means replacing many scalar operations written as loops with array operations that act on whole vectors, matrices, or tensors. The mathematical result may be the same, but the implementation can be much faster because optimized libraries use cache-friendly memory access, parallel instructions, BLAS kernels, and GPU execution.

Vectorization is not merely a speed trick. It also makes shapes explicit. A batch of \(B\) examples should usually be processed as a tensor with a batch axis, not as \(B\) unrelated scalar programs.

\paragraph{Formal setup.}
Let \(X\in\mathbb{R}^{B\times d_{\mathrm{in}}}\) be a batch of inputs, \(W\in\mathbb{R}^{d_{\mathrm{out}}\times d_{\mathrm{in}}}\) be a weight matrix, and \(b\in\mathbb{R}^{d_{\mathrm{out}}}\) be a bias. The batched affine layer is
\[
Y=XW^\top+\mathbf{1}b^\top,
\qquad
Y\in\mathbb{R}^{B\times d_{\mathrm{out}}},
\]
where \(\mathbf{1}\in\mathbb{R}^{B}\) is a vector of ones. Entrywise,
\[
Y_{b,o}=\sum_{i=1}^{d_{\mathrm{in}}}X_{b,i}W_{o,i}+b_o.
\]
The formula keeps the batch axis \(b\), contracts the input feature axis \(i\), and creates an output feature axis \(o\).

Broadcasting is a rule for reusing an array along a compatible size-one or missing axis. In the affine layer, \(b_o\) is broadcast over all batch items. Broadcasting is safe only when the intended repeated axis is clear.

\paragraph{Worked mechanism: vectorized gradients.}
Let \(G=\nabla_Y L\in\mathbb{R}^{B\times d_{\mathrm{out}}}\) be the incoming adjoint for the batched affine layer. The differential is
\[
dY=dXW^\top+X\,dW^\top+\mathbf{1}\,db^\top.
\]
Using \(dL=\langle G,dY\rangle_F\), the gradients are
\[
\boxed{\nabla_X L=GW,\qquad \nabla_W L=G^\top X,\qquad \nabla_b L=G^\top\mathbf{1}.}
\]
Check the shapes:
\[
GW\in\mathbb{R}^{B\times d_{\mathrm{in}}},
\quad
G^\top X\in\mathbb{R}^{d_{\mathrm{out}}\times d_{\mathrm{in}}},
\quad
G^\top\mathbf{1}\in\mathbb{R}^{d_{\mathrm{out}}}.
\]
The vectorized formula is the sum of per-example gradients, computed as matrix multiplications.

\paragraph{Examples and connections.}
A numerical example with \(B=3\), \(d_{\mathrm{in}}=2\), and \(d_{\mathrm{out}}=4\) has
\[
X\in\mathbb{R}^{3\times 2},\quad W\in\mathbb{R}^{4\times 2},\quad Y\in\mathbb{R}^{3\times 4}.
\]
The scalar loop computes twelve output entries. The vectorized expression \(XW^\top+\mathbf{1}b^\top\) computes the same entries through one matrix multiplication plus a broadcasted addition.

In neuroscience, if \(R\in\mathbb{R}^{B\times T\times N}\) stores trials, times, and neurons, and \(D\in\mathbb{R}^{K\times N}\) stores \(K\) decoders, then
\[
Z_{b,t,k}=\sum_{n=1}^N R_{b,t,n}D_{k,n}
\]
produces \(K\) latent readouts for every trial and time. This is a tensor contraction, not \(BTK\) separate hand-written dot products.

In AI, vectorized attention computes many query-key dot products as matrix products. Vectorization also makes automatic differentiation efficient, because the backward pass reuses the same high-performance kernels.

\paragraph{Common misconceptions and failure modes.}
Vectorized code is not automatically correct. Silent broadcasting can hide a missing axis. A transposed weight matrix can produce a shape that works for the wrong reason if dimensions happen to match. Another failure is to optimize arithmetic count while ignoring memory movement; on modern hardware, reading and writing large tensors can dominate computation. Finally, vectorization can reduce clarity if axis names are lost. Good code and good mathematics keep batch, time, neuron, token, and feature axes explicit.

\paragraph{Exercises.}
\begin{enumerate}[label=\textbf{23.4.\arabic*.}, leftmargin=*]
\item \textbf{Mechanical.} For \(X\in\mathbb{R}^{10\times 5}\), \(W\in\mathbb{R}^{3\times 5}\), and \(b\in\mathbb{R}^3\), state the shape of \(Y=XW^\top+\mathbf{1}b^\top\).
\item \textbf{Proof or derivation.} Derive \(\nabla_W L=G^\top X\) for the batched affine layer from \(dL=\langle G,dY\rangle_F\).
\item \textbf{Numerical/computational.} Suppose a scalar loop computes \(Y_{b,o}=\sum_i X_{b,i}W_{o,i}\) for \(B=64\), \(d_{\mathrm{in}}=128\), \(d_{\mathrm{out}}=256\). How many multiply-add terms are evaluated? Which matrix product expresses the same computation?
\item \textbf{Interpretation.} In the neural readout \(Z_{b,t,k}=\sum_n R_{b,t,n}D_{k,n}\), what axes are preserved and what axis is contracted?
\item \textbf{Misconception/debugging.} A program adds a bias vector of shape \(B\) to an activation matrix of shape \(B\times d\), but the intended bias was per feature. Explain the broadcasting mistake and give the correct bias shape.
\end{enumerate}

\paragraph{Closing bridge.}
Automatic differentiation makes derivatives operational, but it does not make numerical computation automatically reliable. The next chapter studies the matrix computations underneath scientific and machine-learning systems: conditioning, stable factorizations, large linear solves, discretization, and error. Those tools explain why some differentiable programs are easy to train and simulate while others are numerically fragile.

\clearpage
